\section{The medium: contracts and clearing}
\label{sec:medium}

\subsection{The debt contract}

The unit is a debt contract: an obligation of a fixed amount from a debtor to a
creditor, with a maturity, in the community's denomination, created by mutual signature
at the point of sale --- the buyer becomes debtor, the seller creditor. It is not a bearer
instrument: a creditor cannot sell, endorse or assign it, which keeps a stake conferred by
one creditor from being resold to another.

Nor is it frozen. Four things move, each as somebody's discharge rather than anybody's
option.

\begin{itemize}[leftmargin=1.4em]
  \item The \textbf{debtor} moves, by transfer or through the support cascade, under the
        consent rule below.
  \item The \textbf{creditor} moves in exactly one circumstance: default on an insured
        obligation, where the claim subrogates to the underwriters whose supply carried it
        (\S\ref{sec:substitution}). Nothing is assigned --- the substitution has already
        made the creditor whole in the same event, and the recovery moves to whoever paid.
  \item The \textbf{maturity} moves only by \emph{extension}, signed by debtor and creditor
        both, because deferring a debt is the creditor's concession to make. A claim the
        cascade routes onto a buyer inherits the earlier of its own date and the sale's, so
        nothing in the alphabet defers a debt without its creditor; an extension past the
        insured horizon of the claim's acceptance keeps the debt and drops the insurance
        (\S\ref{sec:recourse}), on the same two signatures.
  \item The \textbf{outstanding amount}, the \textbf{status} and \textbf{the flow the
        obligation holds} move together as it is discharged, cured, defaulted or routed
        onward, because an insured obligation owes exactly what it holds and the two are
        one fact (\S\ref{sec:verification}).
\end{itemize}

A contract also records whether it is \emph{insured}: a fact about what it holds, not a
label --- it is insured exactly while it holds a reservation
(Definition~\ref{def:reservation}). Decided at acceptance, it can only fall afterwards:
when a downward re-denomination rounds its arcs below one minor unit
(Theorem~\ref{thm:redenominate}), or in the one place a hold is asked for afresh, the
remainder a sale leaves with its seller, if the arcs cannot carry it. A partial payment
shrinks the hold by exactly the share paid, on the arcs that held it, and never re-solves
it. An obligation gains the community's backing only as a \emph{successor} row somebody
signed for, never by the old one changing tier.

\subsection{Origination}

Acceptance takes both parties' signatures and passes one gate: is the amount within the
debtor's capacity, at a maturity inside the insured horizon of its acceptance
(\S\ref{sec:recourse})? If so the obligation is insured: accepting computes an augmenting
flow of exactly the amount into the debtor and holds it (Definition~\ref{def:reservation}),
so the standing that justified this obligation cannot simultaneously justify another. If
not, the creditor may still accept, and the obligation is uninsured: it holds no
reservation and the recourse machinery of \S\ref{sec:recourse} does not stand behind it.
The client says which at the point of decision; a creditor is never uninsured by
accident.

\subsection{What a sale does, in order}
\label{sec:sale}

A sale is a discharge first and a claim second: it creates a claim only out of what it
could not discharge, and the order is the substance of the medium.

\begin{enumerate}[leftmargin=1.6em]
  \item \textbf{Netting.} Whatever the seller already owes the \emph{buyer} extinguishes
        first, oldest first --- settled, or cured if it had matured. Mutual obligations
        cancel rather than route, and as an ordinary discharge by the debtor this applies
        the stake rule of Definition~\ref{def:stake} exactly as any settlement does.
  \item \textbf{The cascade.} What netting did not absorb clears the seller's own
        obligations to third parties, and those of the beneficiaries they have listed, by
        routing them onto the buyer: a debtor swap, which writes no stake in either
        direction (Remark~\ref{rem:only-settlement}).
  \item \textbf{The residue.} Only what neither step absorbed becomes a fresh obligation
        from buyer to seller.
\end{enumerate}

A member who owes is therefore paid by \emph{ceasing to owe}, and a debt is paid here in
the usual course by selling to one's creditor: the debt is extinguished and the stake is
written, so one act discharges the obligation and raises the seller's standing.
\S\ref{sec:adoption} takes up what follows, which is most of the answer to why anybody
accepts this medium.

\subsection{The support cascade}

A member may list \emph{beneficiaries}: other members toward whom their own discharge is
routed. When a supporter delivers value, the discharge waterfills across the beneficiaries
they have listed, in proportion to declared weights, each edge capped at $\kConstNuDrain$
times what that pair has staked in one another (Definition~\ref{def:stake}) and the
recursion bounded at $\kConstMaxCascadeDepth$. The cap is a \emph{stake} and not a settled
volume, for Remark~\ref{rem:not-volume}'s reason: volume is free to fabricate, so a
counter of it would let two cooperating accounts open the drain as wide as they pleased.
It has no floor, so a pair with no settled history drains nothing --- zero is absorbing
here as everywhere else, and escaped the same way, by trading.

This is how a community carries a member through a bad season without anybody making a
gift: the supporter's own production discharges the beneficiary's obligations. What the
supporter gives up is exactly one thing --- the proceeds of their own sale, which would
otherwise have cleared their own debts or become a fresh claim on the buyer --- and
nothing else of theirs is touched, then or later. It carries a member without rebuilding
one: a routed discharge writes no stake in either direction
(Remark~\ref{rem:only-settlement}), so a member cleared this way ends the bad season no
better backed than they entered it --- which is the whole content of the beneficiary's
declaration below, and why it is theirs alone to make.

The cascade records nothing about who contributed what. A co-signer map on each successor
contract would be written, cleared on subrogation, rescaled at every re-denomination,
checked and committed to the state root while no transition reads it, and ledger state
carries what must be enforced; whose production cleared which debt is a question for a
client and a social layer, not a field every node replicates forever.

\begin{property}[A routed claim is never downgraded]
\label{prop:no-downgrade}
Routing an obligation onto the buyer is a \emph{debtor swap}, so it owes the same consent
rule a transfer does. The affected creditors are discovered inside the transition rather
than named by it, and a creditor asked to consent to losing their own recourse would
refuse every time, handing every creditor a veto over their debtor's ordinary sales. The
rule is therefore enforced as a bound on the amount: \textbf{a claim moves only as far as
the buyer can carry it insured}, whatever it was before; the part that does not move stays
with its original debtor, and the budget it did not absorb falls through to the seller as
ordinary new debt. The claim that moves keeps \textbf{the earlier of its own date and the
sale's}: a debtor swap owes the maturity rule a transfer does, and a claim re-dated to the
sale would let two members pass a debt back and forth and push it out of reach for ever,
insured throughout, so that no default ever fires.
\end{property}

Both directions of the bound matter. A claim the buyer can insure is \emph{upgraded} if it
was uninsured before, so no signature is owed for it; a claim the buyer cannot insure does
not travel at all, which stops a member from shedding an obligation onto a freshly minted
key --- two signatures and no delivery --- and walking away with their own standing
restored.

\begin{property}[Beneficiary consent by moderation]
\label{prop:benef}
No member's production is ever routed toward a beneficiary they have not listed, and no
member's obligations are ever cleared by a supporter they have not approved. The two
declarations sit on opposite members --- listing beneficiaries is the supporter's,
approving supporters the beneficiary's --- both are required, and either may be withdrawn.
\end{property}

The two halves answer different questions.

\paragraph{Why the supporter must list.} Without their declaration, approval alone would
be the attack: naming a wealthy member as one's \emph{supporter} would route their
production toward one's own creditors unbidden. Production is directed by the member who
produced it and by nobody else.

\paragraph{Why the beneficiary must approve.} What a drain costs them is the
\emph{standing the obligation would have conferred}. It cannot cost them money, and it is
not what keeps strangers out --- the cap above does that, at zero for a pair that has never
traded whether they approved or not. Routing is a debtor swap the creditor never signs, so
by Remark~\ref{rem:only-settlement} it writes no edge: the debt is discharged and the
member is exactly as backed afterwards as before, while honouring the same obligation
writes that creditor's stake in full. Measured, three loans of $60$ cleared by a supporter
leave a member at the capacity they started with, where the same three honoured take them
from $200$ to $380$. Capacity is a ceiling and what a drain moves is the room beneath it:
the cleared obligation stops reserving what it drew, the release is exact
(Remark~\ref{rem:charge-supply}), and the member ends with the room they had before they
borrowed and the backing they had before the drain --- the $200$ above. The freed flow is
neither awarded to anybody nor lost: the stake edges and the underwriter supply it
occupied are available again, to the successor this sale creates on the buyer, to the
member's own next purchase, or to any other account those arcs reach, first-come as all
capacity is (Remark~\ref{rem:order}).

The beneficiary's declaration is therefore the member choosing \textbf{relief now against
standing later} --- a trade only they can weigh, and one no cap can bound, since the drain
cap limits the amount and not the substitution. One smaller thing rides on it: a
beneficiary's unused share recurses into their own listing, so accepting a supporter
admits that supporter's production into relationships the beneficiary built for their own
use.

\subsection{Settlement, transfer, and conservation}

A contract settles when its outstanding amount reaches zero, by delivery or by mutual
discharge. Settlement releases the reservation the obligation held and applies the stake
rule of Definition~\ref{def:stake}. A partial payment is at least a
$\kConstMaxInstallments$th of the original unless it closes the row, so an obligation is
discharged in at most that many payments and the closing one: a discharge is free, and
free writes are bounded by what somebody paid for (\S\ref{sec:implementation}).

\begin{remark}[Settlement is the only way the graph grows]
\label{rem:only-settlement}
No other transition writes a stake edge. This is what carries
Assumption~\ref{ass:free-signatures}: every quantity that gates participation --- what a
member may confer, their operation-bond headroom, the drain cap of the support cascade,
the floor on proposing --- is read off this graph, so an edge written anywhere but at
settlement is standing nobody earned.
\end{remark}

\emph{Transfer} moves the debtor of a claim: the obligation is assumed by another member,
with their signature, subject to their own capacity gate exactly as at origination. The
creditor does not change, which is why a claim is not a tradable instrument and a stake
conferred by one creditor cannot be resold to another.

A transfer is a discharge of the outgoing debtor, so it must be authorised by whoever
loses if it is wrong --- the creditor, but only when the claim stops being one the
community stands behind. Hence the consent rule: \textbf{the creditor signs a transfer
exactly when the successor would be uninsured.} Where the successor can carry it insured,
the community's recourse follows the claim to its new debtor, and who that debtor is has
become the community's question rather than the creditor's --- in this model's terms not a
change to consent to. Which case applies is a fact about the graph, not a choice anybody
makes.

A transfer also stakes \emph{nothing}: the outgoing debtor did not pay, so
Definition~\ref{def:stake} does not fire, and the graph grows when the successor settles,
against the debtor who actually delivered. Otherwise a ring of members passing one
accepted claim between them would each collect a stake from a creditor who signed once,
and Remark~\ref{rem:only-settlement} would be false.

\begin{invariant}[Conservation]
\label{inv:conservation}
At every height, the sum of outstanding amounts over all active contracts equals the sum
of cached outstanding debt over all accounts, and no transition creates or destroys
obligation except origination (which creates exactly one) and discharge (which reduces
exactly one).
\end{invariant}

\begin{invariant}[Reservation conservation]
\label{inv:res-conservation}
At every height, $\sum \res(a,b)$ equals the total outstanding amount of insured
contracts. A reservation exists if and only if there is a live insured obligation behind
it.
\end{invariant}

Invariant~\ref{inv:res-conservation} is what makes Theorem~\ref{thm:cut} enforced rather
than argued: if reservations could drift from obligations, the cut would bound a quantity
nobody was tracking.

\subsection{Default and cure}

An obligation past maturity with a positive outstanding amount is in \emph{default}, and
the consequences are asymmetric. For an \textbf{insured} obligation the community's
recourse machinery applies (\S\ref{sec:recourse}): \emph{substitution} always, and
arbitration if the parties consented to it. For an \textbf{uninsured} one none of it
applies: the debt stands, the claim stands, and the creditor bears the loss alone
(Remark~\ref{rem:uninsured-alone}).

\emph{Cure} is late discharge: a defaulted obligation may be settled afterwards, which
clears the default and applies the stake rule as any settlement does. Default is a state,
not a verdict; a member who repays late is not permanently marked, and what marks them is
that their counterparties may stake less next time, which the community decides and the
protocol does not impose.

\subsection{The capacity gate}

\begin{invariant}[Capacity]
\label{inv:capacity}
No transition leaves any account set $S$ holding insured obligations whose total exceeds
$\Capa(S)$.
\end{invariant}

This is the only credit limit in the protocol: no per-relationship cap, no per-epoch
volume ceiling, no channel with its own floor and no separate rule for newcomers, because
Definition~\ref{def:capacity} already gives a newcomer zero and Corollary~\ref{cor:sybil}
gives a thousand newcomers zero.
